\documentclass[11pt]{article}
\usepackage[margin=1in]{geometry}
\usepackage{amsmath,amssymb,booktabs,hyperref}
\title{Identifiability of System State and Evidence Quality}
\author{Guillaume Harbonnier}
\date{ED-POMDP Research Edition v1.1.0}
\begin{document}
\maketitle
\begin{abstract}
This note examines when latent system condition $S$ and latent evidence-production quality $E$ can be separated from assurance observations. It provides a non-identifiability counterexample, a heterogeneous-evidence case, and a controlled-intervention case with an explicit intervention noise floor. It supports claim \texttt{CLM-IDENT-001} in the canonical claim registry and does not claim empirical identifiability.
\end{abstract}

\section{Problem}
Let observations follow $P(O\mid S,E)$. The factorization is useful only if distinct values of $S$ and $E$ can be distinguished under stated assumptions.

\section{Case A: non-identifiability}
If two latent pairs induce the same observation law,
\[
P(O\mid S=s_1,E=e_1)=P(O\mid S=s_2,E=e_2),
\]
then repeated evidence from that unchanged channel cannot separate them. More samples reduce variance but do not repair structural aliasing.

\section{Case B: heterogeneous evidence}
Separation can become possible when channels have different sensitivity profiles. A functional test may be primarily sensitive to $S$, while an environment validation probe is primarily sensitive to $E$. The joint observation map must have distinct rows, at least up to the equivalence classes relevant to the decision.

\section{Case C: controlled intervention}
Changing an evidence-production condition while holding the build nominally fixed can increase information about $E$. Yet the system is never perfectly invariant between runs. Flakiness, timing, hidden state, load and environmental drift create an intervention noise floor $\epsilon_{\mathrm{drift}}$.

Hence an intervention identifies $E$ only relative to an explicit stability assumption and only above that noise floor. Paired runs, randomized order, invariant checks and nuisance-variable logging are required before attributing an observation change to $E$.

\section{Practical implications}
Identifiability comes from experimental structure rather than evidence volume alone. Assurance designs should therefore include environment smoke tests, observability and oracle validation, heterogeneous channels, paired cross-environment execution and controlled interventions with nuisance logging.

\section{Current epistemic status}
The $S/E$ factorization is formally meaningful, but empirical identifiability is not established. Step 2 must test identifiable, weakly identifiable and non-identifiable regimes separately.
\end{document}
